\documentclass{article}

% Language setting
% Replace `english' with e.g. `spanish' to change the document language
\usepackage[english]{babel}

% Set page size and margins
% Replace `letterpaper' with `a4paper' for UK/EU standard size
\usepackage[letterpaper,top=2cm,bottom=2cm,left=3cm,right=3cm,marginparwidth=1.75cm]{geometry}
\usepackage[natbib=true, backend=biber, style=authoryear]{biblatex}
\addbibresource{references.bib}

% Useful packages

\usepackage{amsmath}
\usepackage{amsfonts}
\usepackage{graphicx}
\usepackage[colorlinks=true, allcolors=blue]{hyperref}
\usepackage{bm}
\usepackage{csquotes}



\newcommand{\cf}{cf.\xspace}
\newcommand{\eg}{e.g.\xspace}
\newcommand{\ie}{i.e.\xspace}


\title{Assignment\\Neurosymbolic AI -- Semantic Loss\\Solution}
\author{Your Name(s)}
\date{}





\usepackage{minted}


\begin{document}
\maketitle




\section*{Submission Intructions}
\begin{itemize}
	\item For solving this problem you are allowed to work in pairs.
	\item In order to submit your solution use this LaTex template.
	\item Clearly indicate your name(s).
	\item If your work in pairs. Both of you should upload a version to Canvas and each
	\item The uploaded solution has to contain both of your names.
\end{itemize}



\section{Training a Neural Classifier}



\section{Replacing Softmax with Sigmoids}




\section{Encouraging Logical Constraints with the Semantic Loss}

\subsection*{Q3.1}

\begin{minted}{python3}
def exactly_one(n):
	# n is the number of Boolean variables
	...


\end{minted}


\subsection*{Q3.2}


\subsection*{Q3.3}

\begin{minted}{python3}
class LogProbSemiring:

	@staticmethod
	def one():
		raise NotImplementedError()

	@staticmethod
	def zero():
		raise NotImplementedError()

	@staticmethod
	def plus(a, b):
		raise NotImplementedError()

	@staticmethod
	def times(a, b):
		raise NotImplementedError()

	@staticmethod
	def value(a):
		raise NotImplementedError()

	@staticmethod
	def negate(a):
		raise NotImplementedError()
	\end{minted}


\subsection*{Q3.4}








\section{Open-Ended Question}

\begin{description}
	\item[Q4] Come up with your own neurosymbolic problem where you would like to encourage a symbolic constraint on the outputs of a neural network.
\end{description}

Hint: Do not make it too complicated! Our naive implementation does not scale that well. For thos more sophisticated techniques are needed \citep{maene2025klay} but that is a topic for another day.


\section{Neurosymbolic AI and Physics Informed Neural Networks}

\begin{description}
	\item[Q5] Read up on \textit{physics informed neural networks} \citep{raissi2019physics} and compare them to the semantic loss. What are differences? What are commonalities?
	      \\
	      \href{https://www.youtube.com/watch?v=-zrY7P2dVC4}{https://www.youtube.com/watch?v=-zrY7P2dVC4}
\end{description}






\printbibliography
\end{document}